\documentclass{amsart}
\usepackage{amsmath,amssymb,amsthm,hyperref}

\newtheorem{theorem}{Theorem}
\newtheorem{lemma}{Lemma}
\newtheorem{corollary}{Corollary}
\newtheorem{proposition}{Proposition}
\theoremstyle{remark}
\newtheorem*{remark}{Remark}

\begin{document}

\title{An independence bound for planar graphs without the Four Color Theorem}
\maketitle

\begin{center}\textbf{Honest status of the target claim}\end{center}

The problem asks to prove that every planar graph on $n$ vertices has an
independent set of size at least $n/4$, \emph{without} invoking the Four Color
Theorem (4CT). I must report at the outset, with full mathematical honesty, that
the bound $n/4$ in this unconditional, 4CT-free form is \textbf{not known to be
provable by elementary means}: the best published lower bound on the independence
ratio of planar graphs that does not use the 4CT is $3/13 = 0.2307\ldots < 1/4$
(Cranston--Rabern, 2016), improving Albertson's $2/9 = 0.2222\ldots$ (1976). The
bound $n/4$ is tight (the icosahedron has $n=12$ and $\alpha = 3 = n/4$), and
obtaining it unconditionally is regarded as comparable in difficulty to the 4CT
itself. Fabricating a proof of $n/4$ would violate the requirement of grounded,
gap-free reasoning, so I do not do so.

Instead I give two fully complete, self-contained results that genuinely avoid the
4CT, and I delineate precisely (Section~\ref{sec:gap}) where every elementary
method stops short of $n/4$ and why that is the genuine obstruction:
\begin{itemize}
\item[(A)] $\alpha(G)\ge n/5$ for \emph{every} planar graph $G$
  (Theorem~\ref{thm:main}), via the Five Color Theorem proved from first
  principles;
\item[(B)] $\alpha(G)\ge n/4$ for every \emph{$3$-degenerate} planar graph
  (Proposition~\ref{prop:3deg}), in particular for every planar graph of girth
  at least $6$ (Corollary~\ref{cor:girth}); this recovers the exact target bound
  on a broad subclass.
\end{itemize}

\bigskip

Throughout, $G=(V,E)$ is a finite simple graph with $n=|V|$ vertices and $m=|E|$
edges, $\alpha(G)$ denotes the maximum size of an independent set, and
$\chi(G)$ denotes the chromatic number. For a vertex $v$, $N(v)$ is its set of
neighbours, $\deg(v)=|N(v)|$, and $N[v]=N(v)\cup\{v\}$. The \emph{girth} of a
graph with a cycle is the length of a shortest cycle. We fix a plane embedding of
the planar graph $G$ when one is needed.

\section{Pigeonhole reduction from colouring to independence}

\begin{lemma}\label{lem:pigeon}
If a graph $G$ on $n$ vertices satisfies $\chi(G)\le k$, then
$\alpha(G)\ge n/k$.
\end{lemma}

\begin{proof}
By definition of $\chi(G)\le k$ there is a proper colouring
$c\colon V\to\{1,\dots,k\}$, i.e.\ $c(u)\ne c(v)$ whenever $uv\in E$. For each
colour $i\in\{1,\dots,k\}$ let $C_i=c^{-1}(i)=\{v\in V: c(v)=i\}$. The sets
$C_1,\dots,C_k$ partition $V$, so $\sum_{i=1}^{k}|C_i|=n$. By the pigeonhole
principle there is an index $j$ with $|C_j|\ge n/k$. Each $C_i$ is independent:
if $u,v\in C_i$ then $c(u)=c(v)=i$, and properness forces $uv\notin E$. Hence
$C_j$ is an independent set with $|C_j|\ge n/k$, giving $\alpha(G)\ge n/k$.
\end{proof}

Thus it suffices to prove $\chi(G)\le 5$ for every planar graph $G$. (The 4CT is
the statement $\chi(G)\le 4$; we deliberately avoid it and prove the weaker,
elementary Five Color Theorem, which yields the $n/5$ bound.)

\section{Euler's formula, edge counts, and low-degree vertices}

\begin{lemma}\label{lem:edges}
If $G$ is a simple planar graph with $n\ge 3$ vertices and $m$ edges, then
$m\le 3n-6$. More generally, if $G$ is a simple planar graph with $n\ge3$
vertices, $m$ edges, and girth at least $g\ge3$ (in particular $G$ contains a
cycle), then
\begin{equation}\label{eq:girth-edges}
   m \;\le\; \frac{g}{g-2}\,(n-2).
\end{equation}
\end{lemma}

\begin{proof}
We first prove \eqref{eq:girth-edges} for connected $G$ that contains a cycle.
Fix a plane embedding and let $f$ denote the number of faces (including the
unbounded one). Euler's formula for connected plane graphs states
\begin{equation}\label{eq:euler}
   n - m + f = 2 .
\end{equation}
Counting incidences between faces and edge-sides, each edge borders at most two
faces, so
\[
   \sum_{\text{faces } \phi}\operatorname{len}(\partial\phi)\;\le\;2m .
\]
Since $G$ is connected, simple, with girth $\ge g$, and contains a cycle, every
face boundary walk has length at least $g$. Indeed, in a connected plane graph
the boundary walk of a bounded face encloses a region, so it must traverse a
cycle of $G$, and the unbounded face's boundary walk likewise traverses the outer
cycle of $G$; either way each boundary walk contains a cycle of $G$, whose length
is $\ge g$ (and lengths $1$ and $2$ are excluded by simplicity). Bridges, when
present, are traversed twice in the boundary walk of the single face they border,
which only lengthens that walk and does not lower the bound. Therefore
\[
   g\,f \;\le\; \sum_{\text{faces }\phi}\operatorname{len}(\partial\phi)\;\le\;2m,
   \qquad\text{so}\qquad f\le \tfrac{2m}{g}.
\]
Substituting into \eqref{eq:euler},
\[
   2 = n-m+f \le n - m + \tfrac{2m}{g} = n - m\Bigl(1-\tfrac{2}{g}\Bigr)
       = n - m\,\tfrac{g-2}{g},
\]
which rearranges to $m\,\tfrac{g-2}{g}\le n-2$, i.e.\
$m\le \tfrac{g}{g-2}(n-2)$.

Taking $g=3$ gives, for connected $G$ with a cycle, $m\le 3(n-2)=3n-6$. If a
connected $G$ has no cycle it is a tree, so $m=n-1\le 3n-6$ for $n\ge3$. For a
disconnected simple planar graph one may add edges between distinct components
(preserving planarity and simplicity, strictly increasing $m$, fixing $n$) until
connected; the maximiser of $m$ for fixed $n$ is thus connected, and the bound
$m\le 3n-6$ for it dominates the original. This proves the first statement
$m\le3n-6$ for all simple planar $G$ with $n\ge3$. (For \eqref{eq:girth-edges}
with $g>3$ we will only invoke it for connected graphs containing a cycle, the
case just established.)
\end{proof}

\begin{lemma}\label{lem:mindeg}
Every simple planar graph $G$ has a vertex of degree at most $5$.
\end{lemma}

\begin{proof}
If $n\le2$ this is immediate ($\deg(v)\le n-1\le1\le5$). Assume $n\ge3$ and
suppose, for contradiction, that every vertex has degree at least $6$. Then by the
handshake identity $2m=\sum_{v}\deg(v)\ge 6n$, so $m\ge 3n$. This contradicts
Lemma~\ref{lem:edges}, which gives $m\le 3n-6<3n$. Hence some vertex has degree
$\le5$.
\end{proof}

\begin{lemma}\label{lem:girth6deg}
Every simple planar graph $H$ of girth at least $6$ has a vertex of degree at
most $2$. Consequently every planar graph of girth at least $6$ is
$2$-degenerate.
\end{lemma}

\begin{proof}
Let $H$ have $n'$ vertices and $m'$ edges. We prove $m' \le \tfrac32(n'-1)$ for
every simple planar $H$ of girth $\ge6$; then the average degree is
$2m'/n' \le 3(n'-1)/n' < 3$, so some vertex has degree at most $2$.

It suffices to prove $m'\le\tfrac32(n'-1)$ for connected $H$: if $H$ has
components $H_1,\dots,H_t$ with $|V(H_i)|=n_i'$ and $|E(H_i)|=m_i'$, then summing
$m_i'\le\tfrac32(n_i'-1)$ over $i$ gives
$m'=\sum_i m_i'\le \tfrac32\sum_i(n_i'-1)=\tfrac32(n'-t)\le\tfrac32(n'-1)$, using
$t\ge1$.

So let $H$ be connected. If $H$ is a tree then $m'=n'-1\le\tfrac32(n'-1)$. If $H$
contains a cycle then $n'\ge6\ge3$ and \eqref{eq:girth-edges} with $g=6$ gives
$m'\le \tfrac{6}{4}(n'-2)=\tfrac32(n'-2)\le\tfrac32(n'-1)$. In all cases
$m'\le\tfrac32(n'-1)$, as claimed.

Finally, every subgraph of a planar graph of girth $\ge6$ is itself planar of
girth $\ge6$ (removing vertices or edges cannot create shorter cycles and cannot
destroy planarity), so each subgraph has a vertex of degree $\le2$; this is the
definition of $2$-degeneracy.
\end{proof}

\section{The Five Color Theorem}

\begin{theorem}[Five Color Theorem]\label{thm:five}
Every planar graph $G$ satisfies $\chi(G)\le 5$.
\end{theorem}

\begin{proof}
We argue by induction on $n=|V(G)|$. If $n\le5$, assign distinct colours from
$\{1,\dots,5\}$ to the at most five vertices; this is a proper $5$-colouring.

\emph{Inductive step.} Let $n\ge6$ and assume every planar graph on fewer than
$n$ vertices is $5$-colourable. Fix a plane embedding of $G$. By
Lemma~\ref{lem:mindeg}, $G$ has a vertex $v$ with $\deg(v)\le5$. The graph $G-v$
is planar (deleting a vertex preserves planarity) and has $n-1$ vertices, so by
induction it has a proper $5$-colouring $c\colon V(G)\setminus\{v\}\to
\{1,\dots,5\}$. We extend $c$ to $v$.

\medskip
\noindent\emph{Case 1: $\deg(v)\le 4$, or the neighbours of $v$ use at most $4$
distinct colours.} Then at least one colour, say $a$, appears on no neighbour of
$v$. Set $c(v)=a$. Since $a$ differs from the colour of every neighbour of $v$
and $c$ was proper on $G-v$, the extension is a proper $5$-colouring of $G$.

\medskip
\noindent\emph{Case 2: $\deg(v)=5$ and the five neighbours of $v$ receive five
\emph{distinct} colours.} Let the neighbours, listed in the cyclic order in which
the edges $vv_i$ leave $v$ in the plane embedding, be $v_1,v_2,v_3,v_4,v_5$, and
relabel colours so that $c(v_i)=i$ for $i=1,\dots,5$.

Consider colours $1$ and $3$. Let $H_{1,3}$ be the subgraph of $G-v$ induced by
all vertices coloured $1$ or $3$, and let $K$ be the connected component of
$H_{1,3}$ containing $v_1$.

\emph{Subcase 2a: $v_3\notin K$.} Recolour, \emph{within $K$ only}, by swapping
colours $1$ and $3$ (each vertex of $K$ coloured $1$ becomes $3$ and vice versa).
Call the new colouring $c'$. We verify $c'$ is a proper colouring of $G-v$. The
only edges whose endpoints might conflict are those with exactly one endpoint in
$K$. Let $xy$ be such an edge with $x\in K$, $y\notin K$. As $x\in K\subseteq
H_{1,3}$, we have $c(x)\in\{1,3\}$. If $y$ had colour $1$ or $3$, then $y\in
H_{1,3}$ and, being adjacent to $x\in K$, $y$ would lie in the same component
$K$, contradicting $y\notin K$. Hence $c(y)\in\{2,4,5\}$. The swap gives
$c'(x)\in\{1,3\}$ and leaves $c'(y)=c(y)\in\{2,4,5\}$, so $c'(x)\ne c'(y)$. Edges
with both endpoints in $K$ keep proper colours (both endpoints swap $1\leftrightarrow3$
consistently), and edges with both endpoints outside $K$ are untouched. Hence
$c'$ is a proper $5$-colouring of $G-v$. Under $c'$, $v_1\in K$ now has colour
$3$, while $v_3\notin K$ keeps colour $3$; thus the neighbours of $v$ now use the
colours $3,2,3,4,5$, so colour $1$ appears on no neighbour of $v$. Set $c'(v)=1$
to obtain a proper $5$-colouring of $G$.

\emph{Subcase 2b: $v_3\in K$.} Then $v_1$ and $v_3$ lie in the same component of
$H_{1,3}$, so there is a path $P_{1,3}$ from $v_1$ to $v_3$ whose internal
vertices are all coloured $1$ or $3$. Adjoining the edges $vv_1$ and $vv_3$ to
$P_{1,3}$ yields a cycle $C$ drawn in the plane. By the Jordan curve theorem,
the closed curve $C$ separates the plane into a bounded interior and an unbounded
exterior. Because the embedding fixes the cyclic order $v_1,v_2,v_3,v_4,v_5$ of
edges around $v$, the neighbour $v_2$ lies in the cyclic sector strictly between
$v_1$ and $v_3$, whereas $v_4$ (and $v_5$) lie in the complementary sector;
consequently $v_2$ and $v_4$ lie in different regions determined by $C$ -- one
strictly inside $C$ and the other strictly outside. (This is exactly where
planarity is used.)

Now consider colours $2$ and $4$: let $H_{2,4}$ be the subgraph of $G-v$ induced
by vertices coloured $2$ or $4$, and $K'$ the component of $H_{2,4}$ containing
$v_2$. A path from $v_2$ to $v_4$ inside $H_{2,4}$ uses only vertices coloured
$2$ or $4$ and never uses $v$; since the curve $C$ consists solely of $v$ and
vertices coloured $1$ or $3$, such a path is disjoint from $C$ and hence cannot
cross from the interior to the exterior of $C$. As $v_2$ and $v_4$ lie on
opposite sides of $C$, no $\{2,4\}$-coloured path joins them, so $v_4\notin K'$.
We now repeat the swap of Subcase~2a with the colour pair $(2,4)$ and component
$K'$: swapping colours $2$ and $4$ on $K'$ produces a proper $5$-colouring $c''$
of $G-v$ (by the identical verification as in Subcase~2a, with $\{1,3\}$ replaced
by $\{2,4\}$) under which $v_2$ becomes colour $4$ while $v_4$ retains colour $4$.
Then colour $2$ appears on no neighbour of $v$, and setting $c''(v)=2$ gives a
proper $5$-colouring of $G$.

In every case $G$ admits a proper $5$-colouring, completing the induction.
\end{proof}

\section{The independence bound (A)}

\begin{theorem}\label{thm:main}
Every planar graph $G$ on $n$ vertices contains an independent set of size at
least $n/5$; that is, $\alpha(G)\ge n/5$. This argument does not use the Four
Color Theorem.
\end{theorem}

\begin{proof}
By Theorem~\ref{thm:five}, $\chi(G)\le5$. Applying Lemma~\ref{lem:pigeon} with
$k=5$ gives $\alpha(G)\ge n/5$. The only colouring input is the Five Color
Theorem, whose proof uses only Euler's formula, the existence of a low-degree
vertex, and Kempe-chain recolouring; the Four Color Theorem is never invoked.
\end{proof}

\section{The exact bound \texorpdfstring{$n/4$}{n/4} on \texorpdfstring{$3$}{3}-degenerate
planar graphs (B)}\label{sec:3deg}

A graph $G$ is \emph{$d$-degenerate} if every (vertex-induced) subgraph of $G$
contains a vertex of degree at most $d$.

\begin{proposition}\label{prop:3deg}
If $G$ is a $3$-degenerate graph on $n$ vertices, then $\alpha(G)\ge n/4$.
\end{proposition}

\begin{proof}
We argue by induction on $n$. If $n=0$ the empty set is independent and
$0\ge 0/4$. If $1\le n\le4$, any single vertex forms an independent set of size
$1\ge n/4$. Now let $n\ge5$ and assume the statement for all $3$-degenerate
graphs on fewer than $n$ vertices. Since $G$ is $3$-degenerate (taking the whole
graph as the subgraph), there is a vertex $v$ with $\deg_G(v)\le3$. Let
$G'=G-N[v]$, the graph obtained by deleting $v$ together with its at most three
neighbours; thus $|V(G')|\ge n-4$. Every subgraph of $G'$ is a subgraph of $G$,
hence contains a vertex of degree $\le3$, so $G'$ is again $3$-degenerate. By the
induction hypothesis $\alpha(G')\ge |V(G')|/4\ge (n-4)/4$. Let $I'$ be a maximum
independent set of $G'$. No vertex of $I'$ is adjacent to $v$ (all neighbours of
$v$ were deleted in forming $G'$), and $v\notin V(G')$, so $I:=I'\cup\{v\}$ is an
independent set of $G$ with $|I|=|I'|+1\ge (n-4)/4+1 = n/4$. Hence
$\alpha(G)\ge n/4$.
\end{proof}

\begin{corollary}\label{cor:girth}
Every planar graph $G$ of girth at least $6$ satisfies $\alpha(G)\ge n/4$.
\end{corollary}

\begin{proof}
By Lemma~\ref{lem:girth6deg}, $G$ is $2$-degenerate, hence $3$-degenerate.
Proposition~\ref{prop:3deg} then gives $\alpha(G)\ge n/4$.
\end{proof}

\section{Where the elementary method stops short of \texorpdfstring{$n/4$}{n/4}}
\label{sec:gap}

For transparency we record exactly why the elementary toolkit yields $n/5$ for
all planar graphs and the exact $n/4$ only on restricted classes.

\medskip
\noindent\textbf{(i) The colouring route caps at $n/5$.}
Lemma~\ref{lem:pigeon} turns a $k$-colouring into an independent set of size
$n/k$. Without the 4CT the best elementary colouring bound is $\chi\le5$
(Theorem~\ref{thm:five}), giving exactly $n/5$. The 4CT would give $\chi\le4$ and
hence $n/4$, but $\chi\le4$ for planar graphs \emph{is} the 4CT, which is
forbidden here.

\medskip
\noindent\textbf{(ii) The greedy/degeneracy route caps at $n/6$ in general.}
A planar graph is $5$-degenerate (every subgraph has a vertex of degree $\le5$ by
Lemma~\ref{lem:mindeg}), which yields only $\alpha\ge n/(5+1)=n/6$ by the standard
greedy elimination $\alpha(G)\ge n/(d+1)$ for $d$-degenerate $G$. The improvement
to $n/4$ in Proposition~\ref{prop:3deg} requires $3$-degeneracy, which not every
planar graph has.

\medskip
\noindent\textbf{(iii) The obstruction is genuine.}
The recursion of Proposition~\ref{prop:3deg} recovers $n/4$ exactly when a
degree-$\le3$ vertex keeps appearing after each closed-neighbourhood removal.
There exist planar graphs of minimum degree $\ge4$ that are \emph{not}
$3$-degenerate; for them a closed-neighbourhood removal can delete $\ge5$
vertices, destroying the ``one independent vertex per $4$ deleted'' budget. The
icosahedron ($5$-regular, $n=12$, $\alpha=3=n/4$) is the canonical tight
obstruction: being $5$-regular it has no vertex of degree $\le3$, so the
recursion never starts, yet it still meets $n/4$ with equality.

\medskip
\noindent\textbf{(iv) State of the art.}
Closing the gap unconditionally is, to current knowledge, open. The best
published independence ratio for planar graphs avoiding the 4CT is
$3/13=0.2307\ldots<1/4$ (Cranston--Rabern, 2016), improving Albertson's
$2/9=0.2222\ldots$ (1976); both exceed the $1/5$ of Theorem~\ref{thm:main} but
remain below $1/4$. The Albertson--Berman conjecture -- that every planar graph
has an induced forest on $\ge n/2$ vertices, which would yield $\alpha\ge n/4$
because a forest is bipartite -- remains unproven. Accordingly, a fully
unconditional elementary proof of the exact bound $n/4$ for \emph{all} planar
graphs is not currently available, and producing one here would require an
unjustified leap. The strongest complete, gap-free results obtainable by
self-contained means are recorded above: $\alpha(G)\ge n/5$ for all planar $G$
(Theorem~\ref{thm:main}), and $\alpha(G)\ge n/4$ for all $3$-degenerate planar
$G$ (Proposition~\ref{prop:3deg}), in particular for all planar graphs of girth
$\ge6$ (Corollary~\ref{cor:girth}).

\end{document}
